\documentclass[12pt,letterpaper]{article}

% ---- Packages ----
\usepackage[utf8]{inputenc}
\usepackage[T1]{fontenc}
\usepackage{amsmath,amssymb}
\usepackage{newtxtext,newtxmath}  % Times New Roman text + math (replaces mathptmx)
\usepackage[margin=1in]{geometry}
\usepackage{setspace}
\usepackage{titlesec}
\usepackage{enumitem}
\usepackage{hyperref}
\usepackage{microtype}
\usepackage{booktabs}

% ---- Section formatting ----
\titleformat{\section}{\large\bfseries\sffamily}{\thesection.}{0.5em}{}
\titleformat{\subsection}{\normalsize\bfseries\sffamily}{\thesubsection}{0.5em}{}
\titleformat{\subsubsection}{\normalsize\bfseries}{\thesubsubsection}{0.5em}{}

% ---- Spacing ----
\setstretch{1.15}
\setlength{\parskip}{0.5em}
\setlength{\parindent}{0em}

% ---- Custom commands ----
\newcommand{\rhoa}{\rho_{\!A}}
\newcommand{\Ka}{K_{\!A}}
\newcommand{\vin}{v_{\mathrm{in}}}
\newcommand{\rs}{r_s}

\begin{document}

\begin{center}
{\LARGE\bfseries Unifying Theory of Dynamic Aether Mechanics:\\
General Overview}\\[1em]
{\large Erik Otto}
\end{center}

\vspace{1em}

\begin{abstract}
This paper introduces a unified physical framework in which a single dynamic
medium---Aether---with just three free parameters produces all known forces,
special relativity, quantum mechanics, and cosmology.  The framework derives
the Lorentz factor, the magnetic force law, mass--energy equivalence, the
Fresnel drag coefficient, the full gravitational lensing angle, and the
Hawking temperature without fitted parameters, and quantitatively predicts the
five classical experiments historically regarded as disproving aether theories
(Michelson--Morley, Kennedy--Thorndike, Sagnac, stellar aberration, and
Fizeau).  All forces arise from two processes in the medium: \emph{flow}
(producing electromagnetism) and \emph{binding} (producing gravity, the
strong force, and the weak force).  These same processes account for
wave--particle duality, atomic spectra, the Pauli exclusion principle, and
the strong and weak nuclear forces.  At cosmological scales, the Aether's
two-fluid thermal structure provides mechanisms for dark energy, dark matter,
and a quantitative resolution of the Hubble tension.  The theory resolves the
vacuum catastrophe and explains the $10^{42}$ hierarchy between
electromagnetic and gravitational force strengths as a consequence of the
near-perfect energy recycling of the binding cycle.  This paper provides a
high-level overview; detailed derivations, equations, and experimental
analyses are presented in the six companion papers of this series.
\end{abstract}

\vspace{1em}
% ============================================================
\section{Introduction}
% ============================================================

The primary focus of this paper is to describe a theory of dynamic Aether---a fundamental
medium permeating all space---and how its properties may be the underlying cause of
magnetic forces, electric forces, electromagnetic unification, the unification of force
constants, gravity, wave--particle duality, the photoelectric effect, entanglement, energy quantization and atomic spectra, the Pauli exclusion principle, constraints from quantum superfluid structure, topological stability of vortex excitations, the gravity--electromagnetism connection, the strong and weak nuclear forces, compatibility with the Standard Model, mass--energy
equivalence, time dilation, gravitational lensing, the classical aether experiments,
the Fresnel drag coefficient, black holes, Hawking radiation, frame dragging, gravitomagnetic precession, neutron star interiors,
cosmological redshift, dark energy, dark matter, and the anomalous integrated Sachs--Wolfe effect in cosmic voids.
I will also explain the striking similarities between this theory and
special relativity, and should this theory prove correct, these similarities would explain
the success special relativity has had.

Central to this theory are three independent properties of Aether---density, pressure,
and temperature---which are decoupled from one another by the mechanism of Aether binding.
In classical fluids these properties are tightly linked: compress a gas and its density,
pressure, and temperature all rise together.  In Aether they separate because Aether can
undergo a phase transition---binding together to form mass.  When Aether binds, it is
removed from the fluid system entirely: it no longer flows, no longer exerts pressure, and
no longer carries thermal energy.  This binding mechanism is not speculative---it is the
experimentally confirmed process of particle--antiparticle pair formation and annihilation,
reinterpreted as the medium coalescing and dissolving.  The Dynamic Casimir Effect, which
converts virtual particles into real ones through moving boundaries, has been directly
observed \cite{wilson2011}.  The theory proposes that the spinning vortex of each electron
acts as such a moving boundary, driving both transient binding (gravity) and permanent
binding (mass formation).

This paper presents the foundational framework as a high-level overview.
The detailed treatments are developed in six companion papers:
special relativity and the classical aether experiments
(Paper~2~\cite{otto-relativity}),
gravity, gravitational lensing, and black holes
(Paper~3~\cite{otto-gravity}),
magnetic and electric fields
(Paper~4~\cite{otto-em}),
the strong and weak nuclear forces
(Paper~5~\cite{otto-nuclear}),
cosmology including dark energy, dark matter, and the Hubble tension
(Paper~6~\cite{otto-cosmology}),
and quantum mechanics including wave--particle duality, atomic spectra, and
superfluid structure
(Paper~7~\cite{otto-qm}).
Each companion paper contains its own mathematical summary, experimental
support section, and open questions.


% ============================================================
\section{Core Principle: Aether Binding}
\label{sec:core-binding}
% ============================================================

Before describing individual phenomena, the mechanism of Aether binding must be understood,
as it is foundational to the rest of the theory.

Aether exists in two states.  \textbf{Unbound Aether} is a dynamic, flowing medium
that exerts pressure, carries waves, and has measurable density and temperature.
\textbf{Bound Aether} (mass) is Aether that has coalesced into particles; once bound,
it no longer flows or exerts pressure as a fluid, and its energy is locked into the
resulting mass.  This is what we observe as matter.

Like superfluid helium, the unbound Aether has internal structure.
It follows the Landau--Tisza two-fluid model:
a superfluid component (zero viscosity, carrying all particle structure,
electromagnetic fields, and binding physics) coexists with a normal
component (finite viscosity, carrying entropy).  Near matter the
superfluid fraction dominates; in warmer regions the normal fraction
grows.  This two-fluid structure is central to cosmology:
light propagates without dissipation (the superfluid component has zero viscosity,
ensuring transparency), while the normal fraction provides the viscous heating that
drives dark energy.\ Through spin-orbit coupling \cite{volovik2003}, the photon's
microscopic spin-wave character produces macroscopic mass currents indistinguishable
from first-sound perturbations, so the effective propagation speed is the two-fluid
first-sound speed $c_1$---whose outward gradient produces the dark matter effect
(Paper~6~\cite{otto-cosmology}).


\subsection{Three Independent Properties}

The binding mechanism decouples Aether's three fundamental properties
so that each varies independently:

\textbf{Density} is essentially constant on macroscopic scales.  The gravitational
inflow of Aether toward matter does not compress the medium---Bernoulli equilibrium
ensures zero net density change.  The electron's spinning vortex core
(at the healing length scale, approximately 0.2~fm) acts as a rapidly moving
boundary whose circulation velocity drives transient binding---the sub-threshold
Dynamic Casimir Effect mechanism described below.  On macroscopic scales gravity
manifests as an inflow velocity field, not a density gradient.  The speed of light
is therefore constant in gravitational fields.

\textbf{Pressure} separates into two types.  The mechanical pressure (set by the
equation of state---the logarithmic relationship between density and pressure
that is unique in producing a constant bulk modulus) is constant wherever
density is constant.  The thermal pressure
varies spatially: lowest near matter, where binding locks thermal energy into mass,
and highest in voids, where accumulated stellar radiation and viscous dissipation heat
the medium.  Only the thermal pressure varies on cosmological scales, driving the
differential expansion observed as dark energy.

\textbf{Temperature} is lowest near matter, where thermal energy is continuously
locked into mass through permanent binding.  It is highest in deep space voids, where
radiated energy from stars accumulates with no mass formation to absorb it.  Despite
their low baryon density, the high temperature of voids produces high thermal pressure,
which drives cosmic expansion.


\subsection{Transient and Permanent Binding}

Aether binding is driven by the Dynamic Casimir Effect \cite{wilson2011}---the
experimentally confirmed process by which a rapidly moving boundary converts
vacuum fluctuations into real particle pairs.  The spinning vortex core of each
electron acts as such a boundary, with circulation velocities near the healing
length that are sufficient to enhance vacuum fluctuations into transient
particle--antiparticle pairs.

\textbf{Transient binding} is the rapid, continuous cycling of Aether in the
near-core region surrounding stable particles.  When the vortex boundary velocity
is sufficient, nearby Aether momentarily coalesces into a particle--antiparticle
pair---a virtual pair in quantum field theory language.  This pair is unstable
and almost immediately annihilates, restoring the consumed Aether.  The cycle
repeats continuously, producing an oscillating pressure wave (normal--low--normal)
that propagates outward.  At macroscopic scales this oscillation is
indistinguishable from a static attractive field: gravity
(Paper~3~\cite{otto-gravity}).

\textbf{Permanent binding} occurs when the vortex boundary velocity exceeds
the DCE permanence threshold---the point at which the moving boundary pumps
enough energy into a virtual pair to prevent its annihilation.  The result is
stable mass.  This threshold divides the two modes: below it, binding is
transient and produces gravity; above it, binding becomes permanent and produces
matter.  Both modes are the same DCE mechanism at different intensities.

Three experimentally confirmed phenomena directly demonstrate transient binding:
the Casimir effect \cite{casimir1948} (measurable force from virtual pairs between
plates), the Lamb shift \cite{lamb1947} (energy-level shifts from virtual pair
interactions), and the anomalous magnetic moment of the electron (corrections from
virtual pair loops confirmed to thirteen significant digits).  The Schwinger effect
\cite{schwinger1951} and the Unruh effect \cite{unruh1976} provide additional
threshold constraints, and in the Aether model all three---Schwinger, Unruh, and
Hawking radiation---are manifestations of a single binding mechanism
(Paper~3~\cite{otto-gravity}).


% ============================================================
\section{Unification of Force Constants}
\label{sec:force-constants}
% ============================================================

With the binding mechanism established, a central result follows: all fundamental
force constants are properties of a single medium.

In standard physics, the speed of light ($c$), the Coulomb constant ($k$), and the
gravitational constant ($G$) are independent and unrelated.  In the Aether model, all
three are properties of a single medium, and only two are truly independent at the
fundamental level.

The speed of light is the wave speed of the medium, determined by its bulk modulus
(stiffness) and density---identical in form to the speed of sound in a classical medium.
The bulk modulus is an intrinsic property of Aether that does not vary anywhere,
which is experimentally confirmed by the exact linearity of the Lorentz factor and
the Fresnel drag coefficient.

The Coulomb constant measures how strongly the medium responds to radial flow
(the flow vector).  It is proportional to the Aether density times the square of
the flow rate per unit charge---a direct measure of the medium's ``stiffness'' to
radial displacement.

The magnetic permeability is not an independent constant.  It is derived from the
Coulomb constant and the speed of light, because magnetism is the angular structure
that develops when a flow vector moves through the medium.  The magnetic force between
two moving charges equals the electric force scaled by the velocity ratio $v_1 v_2/c^2$.
This was first confirmed experimentally in 1855 by Weber and Kohlrausch
\cite{weber1856}, whose measurement of the ratio of electrostatic to electrodynamic
charge units yielded the speed of light.

The gravitational constant measures the residual asymmetry pressure of the bind--unbind
cycle.  Where the Coulomb constant measures the full pressure of direct flow, the
gravitational constant measures only the tiny fraction of binding energy that escapes
after each cycle nearly perfectly cancels itself.  The enormous ratio of electric to
gravitational force strength (approximately $4.17 \times 10^{42}$ for electrons)
decomposes into the near-perfect cancellation fraction and the ratio of flow to
binding coupling---both calculable in principle from Aether parameters.

The factor $v^2/c^2$ appears in four experimentally confirmed phenomena: time dilation,
the magnetic force, relativistic energy, and gravitomagnetism.  In the Aether model all
four share a common origin: the ratio of kinetic energy density (from motion through
the medium) to elastic energy density (the medium's capacity to absorb disturbances).
This universality is a direct consequence of all forces arising from the same medium.

The fine structure constant, the gravitational coupling constant, and the Planck mass
all receive physical interpretations as ratios of Aether properties.  The fine structure
constant measures the strength of a single electron's flow relative to the quantum scale
of the medium; the gravitational coupling constant measures the strength of the residual
binding pressure at the same scale; and the Planck mass is the scale at which the binding
asymmetry becomes comparable to full quantum disturbances.
The detailed equations and derivations are presented in
Papers~2~\cite{otto-relativity}, 3~\cite{otto-gravity},
and 4~\cite{otto-em}.


% ============================================================
\section{Summary of Individual Papers}
% ============================================================

The following sections summarize the key ideas and results of each companion paper.


\subsection{Paper 2: Special Relativity}

The Aether model reproduces the full mathematical content of special relativity from
the physics of motion through a medium (Paper~2~\cite{otto-relativity}).  When an
object moves through Aether at velocity $v$, its internal waves---which propagate at
the medium's wave speed $c$---must travel longer geometric paths, slowing all internal
processes by the Lorentz factor.  This is not an approximation: the exact Lorentz
factor follows from the Pythagorean theorem applied to wave paths.  Length contraction
arises from the same geometry.  Mass--energy equivalence follows from the work
integral against the velocity-dependent effective inertia.

The theory makes a powerful prediction: because the Aether's bulk modulus is constant,
the Lorentz factor must be exactly linear in $v^2/c^2$ with no higher-order corrections.
This has been confirmed from GPS velocities to 0.9994$c$ (CERN muons) with no
deviation detected to parts in $10^8$.  The Fresnel drag coefficient extends this
confirmation across material density ratios up to nearly 6:1.

The five classical experiments historically regarded as disproving aether theories---Michelson--Morley,
Kennedy--Thorndike, Sagnac, stellar aberration, and Fizeau---are all quantitatively
predicted by the model.  The null results follow because the Aether's constant bulk
modulus produces exactly the Lorentz transformations, making the local rest frame
undetectable.


\subsection{Paper 3: Gravity}

Gravity arises from the tiny residual asymmetry of the transient bind--unbind cycle
(Paper~3~\cite{otto-gravity}).  Each cycle nearly perfectly recycles its energy, but a
small fraction propagates outward as a low-pressure wave.  At macroscopic scales the
high cycling frequency makes this indistinguishable from a static attractive field.

Mass drives a steady-state Aether inflow at escape velocity.  A static observer
sitting in this flowing medium experiences time dilation from the relative motion
between their wave-based structure and the medium---gravitational time dilation is
velocity time dilation.  A freely falling observer co-moves with the Aether and
experiences no time dilation, recovering the equivalence principle.  The event horizon
of a black hole is the sonic horizon where the inflow reaches the wave speed.

From this acoustic metric (Painlev\'e--Gullstrand geometry), the theory derives the
full gravitational lensing angle, Hawking radiation as asymmetric transient binding
at the horizon, the photon sphere radius matching Event Horizon Telescope observations,
and frame dragging as aggregate vortex co-rotation (confirmed by Gravity Probe~B
\cite{everitt2011}).

The theory resolves the vacuum catastrophe: quantum field theory's enormous vacuum
energy density is reinterpreted as the Aether's bulk modulus, a real and measurable
quantity that determines the speed of light rather than curving spacetime.


\subsection{Paper 4: Magnetic and Electric Fields}

The electric field is identified as the Aether flow vector---the direction and magnitude
of Aether flow driven by charged particles (Paper~4~\cite{otto-em}).  The magnetic
field is the flow angle---the rotational geometry (vorticity) that develops when a
flow vector moves through the medium.  The relationship between them is determined
by the ratio of source velocity to wave speed: the magnetic force equals the electric
force scaled by $v_1 v_2/c^2$.  The magnetic permeability is therefore not an
independent constant but a derived quantity.

This identification explains why magnetic field lines are always closed
(the mathematical identity that the divergence of a curl is zero), why monopoles
have never been found, and why electric and magnetic fields mix under changes of
reference frame.  The theory accounts for ferromagnetism through Coulomb energy
mediated by Pauli exclusion, the Curie temperature through thermal disruption
of vortex alignment, paramagnetism, diamagnetism, and the Meissner effect through
Cooper pair cancellation of net Aether flow.


\subsection{Paper 5: Strong and Weak Nuclear Forces}

The strong force is carried by compressed Aether bridges---physical strands of
extremely compressed superfluid connecting quark vortex endpoints
(Paper~5~\cite{otto-nuclear}).  These bridges are compressed by a factor of
approximately $7 \times 10^7$ relative to ambient density, and their energy
per unit length (string tension) matches the measured value from lattice QCD.
Confinement follows because pulling quarks apart stretches the bridge until it
snaps, creating a new quark--antiquark pair rather than free quarks.  Asymptotic
freedom has a geometric explanation: the bridges are ``all core'' (their diameter
is comparable to the healing length), so at separations smaller than the bridge
diameter the condensate perturbation has not fully formed and coupling is weak.

The theory reproduces the Lüscher term (the universal correction to the confining
potential from transverse zero-point energy of bridge oscillations), the flux tube
width measured by lattice QCD, the string-breaking distance, and the Schwinger
pair-creation rates that explain why light quarks dominate the pion cloud while heavy
quarks are exponentially suppressed.

The weak force arises from vortex topology changes---reconnections of the electron's
internal vortex structure that change particle identity.  Its short range reflects the
energy cost of core rearrangement, set by the W and Z boson masses.  Color charge
corresponds to the three independent spatial orientations of a bridge.

All four forces are unified as different manifestations of Aether mechanics:
electromagnetic forces from flow-vector pressure, magnetic forces from flow-angle
pressure, the strong force from compressed Aether bridges, the weak force from
vortex topology transitions, and gravity from the bind--unbind asymmetry.


\subsection{Paper 6: Cosmology}

The cosmological framework addresses black holes, neutron stars, and large-scale
structure (Paper~6~\cite{otto-cosmology}).

\textbf{Black holes.}  The event horizon is a sonic horizon where the Aether inflow
reaches the wave speed.  Hawking radiation arises from asymmetric transient binding
at this boundary.  The photon sphere and black hole shadow match Event Horizon
Telescope observations of M87* and Sgr~A*.

\textbf{Neutron stars} serve as Aether laboratories: their extreme densities
(approaching the bridge compression ratio) test the equation of state in regimes
inaccessible on Earth.  Pulsar glitches are explained as quantized vortex
avalanches---the same quantized circulation that stabilizes electrons.
The logarithmic equation of state makes specific predictions for tidal
deformability measurable by gravitational wave observations.

\textbf{Cosmological redshift} results from real Doppler shift due to the
thermal expansion of voids, not from changes in medium properties during propagation.
\textbf{Dark energy} is the thermal pressure of hot void Aether: the normal
(viscous) component of the two-fluid model dissipates cosmic-flow energy into heat,
raising void temperatures and driving accelerating expansion.
\textbf{Dark matter} is a two-fluid first-sound correction: the Landau two-fluid
speed of first sound increases outward from galaxies (where the normal fraction
increases with temperature), producing an additional centripetal acceleration that
mimics dark matter halos.  This naturally produces cored profiles.\ Because photons at macroscopic wavelengths
propagate at $c_1$ (via spin-orbit coupling; Paper~7~\cite{otto-qm}), they follow
geodesics of the same acoustic metric as matter, and rotation curves and lensing are
automatically consistent.

The \textbf{Hubble tension} receives a quantitative resolution: temperature-dependent
dark energy in voids produces a local expansion rate approximately 5~km/s/Mpc above
the cosmic average, matching the observed discrepancy.  The Aether framework also
makes the large local void (KBC void) a natural prediction rather than a $6\sigma$
anomaly.  DESI observations of evolving dark energy are consistent with the
thermal pressure model.


\subsection{Paper 7: Quantum Mechanics}

The quantum mechanical framework addresses wave--particle duality, atomic structure,
and the superfluid properties of the medium (Paper~7~\cite{otto-qm}).

\textbf{Photons} are spin-wave excitations of the spinor Aether condensate---transverse
oscillations of the medium's internal spin degree of freedom.  Derrick's theorem
forbids stable localized particles in the logarithmic condensate, so photons propagate
as extended waves in free space.  Detection occurs when the wave enters the
strong-inflow environment near matter, where Dynamic Casimir Effect dynamics
facilitate condensation into a localized quantum at a single atom.

\textbf{Atomic spectra} arise because atoms act as resonant cavities: the Aether inflow
toward the nucleus creates an acoustic potential well that confines electron waves to
standing-wave modes.  The quantization condition reproduces the hydrogen energy levels,
the virial theorem, selection rules, fine and hyperfine structure, and the Zeeman and
Stark effects.

The \textbf{Pauli exclusion principle} follows from vortex topology: electrons are
quantized vortex rings, and two same-spin vortices in the same orbital would create
a doubly-quantized vortex that is dynamically unstable (confirmed in BEC experiments).
Opposite-spin vortices can coexist because their combined flow pattern remains stable.

The \textbf{quantum superfluid structure} of the Aether is characterized by a
logarithmic equation of state (the unique form producing a density-independent bulk
modulus), a healing length of 0.17--0.28~fm (set by the Aether quantum mass of
500--850~MeV/$c^2$), and quantized circulation.  The electron mass is derived from
vortex ring energy with no fitted parameters beyond the two fundamental Aether
constants.  The Schr\"odinger equation emerges as the Madelung
hydrodynamic equations of the condensate, with Wallstrom's quantization objection
resolved by the physical quantization of circulation.


% ============================================================
\section{Compatibility with the Standard Model}
% ============================================================

The Standard Model of particle physics is the most precisely tested scientific
theory in history.  Any theory proposing a physical substrate for quantum phenomena
must reproduce---not contradict---its successful predictions.  The Aether framework
is not an alternative to the Standard Model but a deeper layer beneath it: the
Standard Model becomes the effective field theory describing low-energy excitations
of the Aether condensate.

The connections fall into three categories:
\textbf{Derived (D)}---results that follow mathematically from the Aether's
fundamental parameters through explicit derivation chains;
\textbf{Analogical (A)}---results established in other condensed matter systems
(superfluid helium-3, spin liquids, cold atomic gases) that are expected to carry
over by structural analogy but have not yet been derived for this specific system;
and \textbf{Proposed (P)}---identifications that are physically motivated but
require future quantitative derivation.

\textbf{Quantum field theory as condensate dynamics.}
The Standard Model's field excitations are identified with condensate excitations:
the electron field with the spinor condensate wave function (\textbf{P}; the Dirac
equation has not been derived from condensate dynamics), the photon field with
flow disturbances propagating at the wave speed (\textbf{D}), and virtual particles
with transient binding events (\textbf{P}).  The QFT formalism is retained and
given a physical ontology---this is a reinterpretation, not a re-derivation.

\textbf{Emergence of gauge symmetries.}
The Standard Model's gauge group is postulated as fundamental; in the Aether
framework it emerges from the condensate order parameter.  U(1) electromagnetic
symmetry arises from conservation of circulation in the superfluid (\textbf{A};
demonstrated by Volovik \cite{volovik2003} in helium-3).  SU(2) weak symmetry
arises from the spinor structure of the two-component order parameter (\textbf{A}).
SU(3) color symmetry emerges from three spatial orientations of the binding bridge
(\textbf{P}; the geometric identification is clear but reproducing the non-Abelian
structure remains open).

\textbf{The Higgs mechanism.}
The Higgs field corresponds to the amplitude of the condensate order parameter
(\textbf{A}).  Spontaneous symmetry breaking occurs because the Aether adopts a
definite ground-state density (\textbf{A}).  The Higgs boson is the amplitude mode
of the condensate---the same mode observed in superconductors and cold atomic gases
(\textbf{D}).  Quantitatively matching the 125~GeV Higgs mass remains open (\textbf{P}).

\textbf{Renormalization and the physical cutoff.}
The ultraviolet cutoff is the healing length---the scale below which the continuum
description fails (\textbf{D}).  Renormalization becomes the physically correct
procedure of restricting calculations to valid length scales (\textbf{D}).  The
ultraviolet divergences of standard QFT never arise because the theory is finite
at all scales.

\textbf{Lorentz invariance.}
The Aether has a preferred frame, but length contraction and time dilation make it
undetectable by any local experiment (\textbf{D}).  Lorentz symmetry is emergent
rather than fundamental---precisely the relationship Volovik demonstrated for
quasiparticles in superfluid helium-3 (\textbf{A}).  Deviations could appear at
energies approaching the healing length scale, but current observations
constrain the spin healing length to energy scales above the Planck energy.

\textbf{What the Standard Model leaves unexplained.}
The Standard Model has approximately 19 free parameters.  The Aether framework
proposes that these reduce to consequences of a smaller set of medium properties:
two free Aether parameters (ambient density and quantum mass) plus the
gravitational asymmetry fraction.  The electron mass has been derived from
vortex energy (\textbf{D}); the full particle mass spectrum, mixing angles,
and coupling constants remain open (\textbf{P}).

\textbf{Open challenges.}
Quantitative QCD reproduction (meson and baryon masses from bridge geometry),
parity violation in the weak interaction, the CKM and PMNS mixing matrices,
and CPT symmetry all require further development.  These are areas where the
framework has not yet been developed, not areas where it conflicts with
established physics.  The relationship is analogous to molecular dynamics and
fluid mechanics: the macroscopic theory (the Standard Model) remains valid, while
the microscopic theory (the Aether model) provides the physical origin of its
parameters and symmetries.


% ============================================================
\section*{Note on Mathematical Content}
% ============================================================

The complete mathematical framework---equations, parameter tables, constraint
summaries, and derivation chains---is presented in each companion paper.
Each paper contains its own Mathematical Summary section collecting all
quantitative relationships developed therein.  The fundamental Aether parameters,
equation of state, and key constraints are summarized consistently across all
seven papers.


% ============================================================
\section*{Acknowledgments}
\addcontentsline{toc}{section}{Acknowledgments}
% ============================================================

The author acknowledges the use of Claude (Anthropic) for \LaTeX{} formatting
and editorial assistance.

\bigskip

% ============================================================
% BIBLIOGRAPHY
% ============================================================

\begin{thebibliography}{99}

\bibitem{wilson2011}
Wilson, C.~M., Johansson, G., Pourkabirian, A., et al.,
``Observation of the dynamical Casimir effect in a superconducting circuit,''
\textit{Nature}, \textbf{479}, 376--379 (2011).

\bibitem{casimir1948}
Casimir, H.~B.~G.,
``On the attraction between two perfectly conducting plates,''
\textit{Proc.\ Kon.\ Ned.\ Akad.\ Wetensch.}, \textbf{51}, 793--795 (1948).

\bibitem{lamb1947}
Lamb, W.~E., Jr.\ and Retherford, R.~C.,
``Fine Structure of the Hydrogen Atom by a Microwave Method,''
\textit{Physical Review}, \textbf{72}, 241--243 (1947).

\bibitem{schwinger1951}
Schwinger, J.,
``On Gauge Invariance and Vacuum Polarization,''
\textit{Physical Review}, \textbf{82}, 664--679 (1951).

\bibitem{unruh1976}
Unruh, W.~G.,
``Notes on black-hole evaporation,''
\textit{Physical Review D}, \textbf{14}, 870--892 (1976).

\bibitem{maxwell1865}
Maxwell, J.~C.,
``A Dynamical Theory of the Electromagnetic Field,''
\textit{Phil.\ Trans.\ Royal Society of London}, \textbf{155}, 459--512 (1865).

\bibitem{weber1856}
Weber, W.\ and Kohlrausch, R.,
``Ueber die Elektricit\"{a}tsmenge, welche bei galvanischen Str\"{o}men durch den Querschnitt der Kette fliesst,''
\textit{Annalen der Physik und Chemie}, \textbf{99}, 10--25 (1856).

\bibitem{michelson1887}
Michelson, A.~A.\ and Morley, E.~W.,
``On the Relative Motion of the Earth and the Luminiferous Ether,''
\textit{American Journal of Science}, \textbf{34}, 333--345 (1887).

\bibitem{kennedy1932}
Kennedy, R.~J.\ and Thorndike, E.~M.,
``Experimental Establishment of the Relativity of Time,''
\textit{Physical Review}, \textbf{42}, 400--418 (1932).

\bibitem{sagnac1913}
Sagnac, G.,
``L'\'{e}ther lumineux d\'{e}montr\'{e} par l'effet du vent relatif d'\'{e}ther dans un interf\'{e}rom\`{e}tre en rotation uniforme,''
\textit{Comptes Rendus}, \textbf{157}, 708--710 (1913).

\bibitem{bradley1728}
Bradley, J.,
``A Letter from the Reverend Mr.\ James Bradley \ldots\ Giving an Account of a New Discovered Motion of the Fix'd Stars,''
\textit{Phil.\ Trans.\ Royal Society}, \textbf{35}, 637--661 (1728).

\bibitem{fizeau1851}
Fizeau, H.,
``Sur les hypoth\`{e}ses relatives \`{a} l'\'{e}ther lumineux,''
\textit{Comptes Rendus}, \textbf{33}, 349--355 (1851).

\bibitem{eddington1920}
Dyson, F.~W., Eddington, A.~S., and Davidson, C.,
``A Determination of the Deflection of Light by the Sun's Gravitational Field, from Observations Made at the Total Eclipse of May 29, 1919,''
\textit{Phil.\ Trans.\ Royal Society A}, \textbf{220}, 291--333 (1920).

\bibitem{everitt2011}
Everitt, C.~W.~F., DeBra, D.~B., Parkinson, B.~W., et al.,
``Gravity Probe B: Final Results of a Space Experiment to Test General Relativity,''
\textit{Physical Review Letters}, \textbf{106}, 221101 (2011).

\bibitem{eht2019}
Event Horizon Telescope Collaboration,
``First M87 Event Horizon Telescope Results.\ I.\ The Shadow of the Supermassive Black Hole,''
\textit{The Astrophysical Journal Letters}, \textbf{875}, L1 (2019).

\bibitem{eht2022}
Event Horizon Telescope Collaboration,
``First Sagittarius~A* Event Horizon Telescope Results.\ I.\ The Shadow of the Supermassive Black Hole at the Center of the Milky Way,''
\textit{The Astrophysical Journal Letters}, \textbf{930}, L12 (2022).

\bibitem{ligo2016}
Abbott, B.~P.\ et al.\ (LIGO Scientific Collaboration and Virgo Collaboration),
``Observation of Gravitational Waves from a Binary Black Hole Merger,''
\textit{Physical Review Letters}, \textbf{116}, 061102 (2016).

\bibitem{steinhauer2019}
Mu\~{n}oz de Nova, J.~R., Golubkov, K., Kolobov, V.~I., and Steinhauer, J.,
``Observation of thermal Hawking radiation and its temperature in an analogue black hole,''
\textit{Nature}, \textbf{569}, 688--691 (2019).

\bibitem{lux2017}
Akerib, D.~S.\ et al.\ (LUX Collaboration),
``Results from a search for dark matter in the complete LUX exposure,''
\textit{Physical Review Letters}, \textbf{118}, 021303 (2017).

\bibitem{xenon2018}
Aprile, E.\ et al.\ (XENON Collaboration),
``Dark Matter Search Results from a One Ton-Year Exposure of XENON1T,''
\textit{Physical Review Letters}, \textbf{121}, 111302 (2018).

\bibitem{pandax2017}
Cui, X.\ et al.\ (PandaX-II Collaboration),
``Dark Matter Results From 54-Ton-Day Exposure of PandaX-II Experiment,''
\textit{Physical Review Letters}, \textbf{119}, 181302 (2017).

\bibitem{desi2025}
DESI Collaboration,
``DESI DR2 Results II: Measurements of baryon acoustic oscillations and cosmological constraints,''
\textit{Phys.\ Rev.\ D}, \textbf{112}, 083515 (2025).

\bibitem{kovacs2022}
Kov\'acs, A., et al.,
``The DES view of the Eridanus supervoid and the CMB cold spot,''
\textit{Mon.\ Not.\ R.\ Astron.\ Soc.}, \textbf{510}, 216--229 (2022).

\bibitem{miller2021}
Miller, M.~C., et~al.,
``The radius of PSR J0740+6620 from NICER and XMM-Newton data,''
\textit{The Astrophysical Journal Letters}, \textbf{918}, L28 (2021).

\bibitem{riley2021}
Riley, T.~E., et~al.,
``A NICER view of the massive pulsar PSR J0740+6620
informed by radio timing and XMM-Newton spectroscopy,''
\textit{The Astrophysical Journal Letters}, \textbf{918}, L27 (2021).

\bibitem{ligo2017gw170817}
Abbott, B.~P., et~al.\ (LIGO/Virgo Collaboration),
``GW170817: observation of gravitational waves from a
binary neutron star inspiral,''
\textit{Physical Review Letters}, \textbf{119}, 161101 (2017).

\bibitem{svancara2024}
Svancara, P., Smaniotto, P., Solidoro, L., et al.,
``Rotating curved spacetime signatures from a giant quantum vortex,''
\textit{Nature} (2024). DOI: 10.1038/s41586-024-07176-8.

\bibitem{hathaway2021}
Hathaway, G.~D.\ and Williams, L.~L.,
``Null-result test for effect on weight from large electrostatic charge,''
\textit{Physics}, \textbf{3}, 160--173 (2021).

\bibitem{tajmar2024}
Tajmar, M., K\"ossling, M., and Neunzig, O.,
``In-depth experimental search for a coupling between gravity and electromagnetism with steady fields,''
\textit{Sci.\ Rep.}, \textbf{14}, 18647 (2024).

\bibitem{volovik2003}
Volovik, G.~E.,
\textit{The Universe in a Helium Droplet}
(Oxford University Press, Oxford, 2003).

\bibitem{otto-relativity}
Otto, E.,
``Unifying Theory of Dynamic Aether Mechanics: Special Relativity,''
(2026).

\bibitem{otto-gravity}
Otto, E.,
``Unifying Theory of Dynamic Aether Mechanics: Gravity,''
(2026).

\bibitem{otto-em}
Otto, E.,
``Unifying Theory of Dynamic Aether Mechanics: Magnetic and Electric Fields,''
(2026).

\bibitem{otto-nuclear}
Otto, E.,
``Unifying Theory of Dynamic Aether Mechanics: Strong and Weak Nuclear Forces,''
(2026).

\bibitem{otto-cosmology}
Otto, E.,
``Unifying Theory of Dynamic Aether Mechanics: Cosmology,''
(2026).

\bibitem{otto-qm}
Otto, E.,
``Unifying Theory of Dynamic Aether Mechanics: Quantum Mechanics,''
(2026).

\end{thebibliography}

\end{document}
